\documentclass{article}
\usepackage{graphicx} % Required for inserting images

\title{FundIA TP1}
\author{Luigi Jacob}
\date{September 2026}

\begin{document}

\maketitle

\section{Tarefa 1 -- Modelagem do Problema}

O problema da Ponte e da Tocha pode ser modelado como um problema de busca em um grafo de estados. Seja o grafo ponderado

\[
G = (V,E),
\]

em que $V$ representa o conjunto de estados possíveis do problema e $E$ representa o conjunto de transições válidas entre esses estados. Cada aresta possui um custo associado, correspondente ao tempo necessário para realizar a travessia que produz a transição.

Dessa forma, a solução do problema consiste em encontrar um caminho no grafo que conecte o estado inicial ao estado objetivo. Como as travessias possuem tempos diferentes, os custos associados às arestas também podem ser diferentes.

\subsection{Representação de Estado}

Cada estado deve armazenar a posição das quatro pessoas e da tocha. Para isso, um estado $S$ é representado pela tupla

\[
S = (A,B,C,D,T),
\]

em que $A$, $B$, $C$ e $D$ representam as quatro pessoas e $T$ representa a posição da tocha.

Cada elemento da tupla assume um valor binário:

\begin{itemize}
    \item $0$: margem inicial;
    \item $1$: margem final.
\end{itemize}

Portanto, o conjunto de estados possíveis pode ser definido como

\[
V \subseteq \{0,1\}^{5}.
\]

Por exemplo, o estado

\[
S = (1,1,0,0,1)
\]

indica que as pessoas $A$ e $B$, juntamente com a tocha, encontram-se na margem final da ponte, enquanto $C$ e $D$ permanecem na margem inicial.

É importante observar que os valores binários representam a \textbf{posição atual} de cada elemento, e não simplesmente se uma pessoa já realizou uma travessia. Isso é necessário porque, durante a solução do problema, pessoas podem atravessar a ponte em ambos os sentidos.

\subsection{Estado Inicial e Estado Objetivo}

No início do problema, todas as pessoas e a tocha encontram-se na margem inicial. Assim, o estado inicial é

\[
S_0 = (0,0,0,0,0).
\]

O objetivo consiste em levar todas as quatro pessoas para a margem final. Como toda travessia necessita da tocha, ao término da solução ela também estará nessa margem. Portanto, o estado objetivo é

\[
S_G = (1,1,1,1,1).
\]

No grafo de estados, $S_0$ corresponde ao vértice de origem da busca e $S_G$ ao vértice objetivo.

\subsection{Função Sucessora}

A função sucessora determina, para um determinado estado $S$, quais estados podem ser alcançados por meio de uma única travessia válida.

Seja

\[
S=(A,B,C,D,T).
\]

Primeiramente, identifica-se o conjunto $P_T$ formado pelas pessoas que estão na mesma margem que a tocha:

\[
P_T = \{p \in \{A,B,C,D\} \mid p = T\}.
\]

Uma ação válida consiste em selecionar uma ou duas pessoas pertencentes a $P_T$. Assim, para um conjunto de pessoas $M$ selecionado para realizar a travessia, deve-se satisfazer

\[
M \subseteq P_T
\]

e

\[
1 \leq |M| \leq 2.
\]

A função sucessora gera todas as combinações válidas de uma ou duas pessoas que estejam na mesma margem que a tocha. Cada uma dessas combinações produz um novo estado e, consequentemente, uma aresta no grafo.

Formalmente, pode-se representar a função sucessora como

\[
Succ (S) =
\{(M,S',c(M)) \mid
M \subseteq P_T,\;
1 \leq |M| \leq 2
\},
\]

em que $M$ representa a ação realizada, $S'$ representa o estado resultante e $c(M)$ representa o custo da travessia.

\subsection{Ações e Transições}

Uma ação consiste no deslocamento de uma ou duas pessoas da margem em que se encontra a tocha para a margem oposta.

Quando uma pessoa participa da travessia, sua posição é invertida. Se $x_p$ representa a posição da pessoa $p$, então

\[
x'_p = 1-x_p,
\qquad \forall p \in M.
\]

As pessoas que não participam da travessia permanecem na mesma margem:

\[
x'_p=x_p,
\qquad \forall p \notin M.
\]

Como toda travessia exige a presença da tocha, sua posição também é invertida:

\[
T'=1-T.
\]

Dessa forma, cada ação válida define uma transição entre dois vértices do grafo.

Por exemplo, considerando o estado inicial

\[
(0,0,0,0,0),
\]

caso as pessoas $A$ e $B$ atravessem juntas, ocorre a transição

\[
(0,0,0,0,0)
\xrightarrow{\{A,B\}}
(1,1,0,0,1).
\]

O novo estado constitui outro vértice do grafo e a travessia realizada corresponde à aresta que conecta os dois estados.

\subsection{Custos das Transições}

Cada pessoa possui um tempo específico de travessia:

\[
t(A)=1,\qquad
t(B)=2,\qquad
t(C)=5,\qquad
t(D)=10.
\]

Quando apenas uma pessoa atravessa, o custo da transição corresponde ao seu próprio tempo:

\[
c(\{p\})=t(p).
\]

Quando duas pessoas atravessam juntas, ambas precisam se deslocar na velocidade da pessoa mais lenta. Portanto, o custo da transição é dado por

\[
c(M)=\max_{p\in M} t(p).
\]

Por exemplo,

\[
c(\{A,B\})=\max(1,2)=2,
\]

enquanto

\[
c(\{A,D\})=\max(1,10)=10.
\]

Assim, cada transição do grafo pode ser representada por uma aresta ponderada

\[
S \xrightarrow[M]{c(M)} S',
\]

na qual $S$ é o estado de origem, $S'$ é o estado resultante, $M$ é a ação realizada e $c(M)$ é o peso da aresta.

Consequentemente, uma solução para o problema corresponde a um caminho

\[
S_0 \rightarrow S_1 \rightarrow \cdots \rightarrow S_G
\]

no grafo de estados. O custo total desse caminho é a soma dos custos das travessias realizadas:

\[
C = \sum_{i=0}^{k-1} c(S_i,S_{i+1}).
\]

O objetivo final consiste, portanto, em encontrar um caminho entre $S_0$ e $S_G$, sendo que, para os métodos de busca orientados por custo, busca-se especificamente o caminho cujo custo total seja mínimo.
```

\section{Algoritmos de Busca}

Para solucionar o problema da Ponte e da Tocha, foram
considerados quatro algoritmos de busca: Busca em
Profundidade (DFS), Busca em Largura (BFS), Busca de
Custo Uniforme (UCS) e Busca A*.

Os algoritmos utilizam o mesmo grafo de estados,
diferenciando-se pelo critério adotado para selecionar
o próximo nó a ser explorado.

\subsection{Busca em Profundidade (DFS)}

A Busca em Profundidade, ou Depth-First Search (DFS),
é um algoritmo de busca não informada que explora
um ramo do grafo até encontrar o estado objetivo
ou até não ser possível avançar para novos nós.

Quando não existem sucessores ainda não explorados,
o algoritmo retrocede para investigar outros caminhos,
realizando o processo conhecido como backtracking.

O DFS utiliza uma estrutura de pilha LIFO
(Last In, First Out), na qual o último elemento
inserido é o primeiro a ser removido.

Como não considera os custos das transições para
determinar a ordem de exploração, o DFS não garante
encontrar a solução de menor custo.

Além disso, em grafos que possuem ciclos, é necessário
controlar os estados visitados para evitar que o
algoritmo explore repetidamente os mesmos caminhos.

\subsection{Busca em Largura (BFS)}

A Busca em Largura, ou Breadth-First Search (BFS),
é um algoritmo de busca não informada que explora
os nós do grafo em ordem crescente de profundidade.

Diferentemente do DFS, o BFS explora todos os nós
de um determinado nível antes de avançar para
o nível seguinte.

Para isso, utiliza uma estrutura de fila FIFO
(First In, First Out), na qual o primeiro elemento
inserido é o primeiro a ser removido.

Em um grafo finito, com controle adequado dos
estados visitados, o BFS encontra uma solução,
caso exista um caminho entre o estado inicial
e o estado objetivo.

O BFS garante encontrar uma solução com o menor
número de transições. Entretanto, como os custos
das travessias podem ser diferentes, essa solução
não necessariamente apresenta o menor tempo total.

\subsection{Busca de Custo Uniforme (UCS)}

A Busca de Custo Uniforme, ou Uniform-Cost Search
(UCS), é um algoritmo de busca não informada que
considera os custos das transições para determinar
a ordem de exploração dos nós.

Diferentemente do BFS, que prioriza os nós de menor
profundidade, o UCS seleciona o nó com o menor
custo acumulado desde o estado inicial.

Esse custo é representado pela função

\[
g(n)=\sum_{i=0}^{k-1}c(S_i,S_{i+1}),
\]

em que $g(n)$ corresponde ao custo acumulado
para alcançar o nó $n$ e $c(S_i,S_{i+1})$
representa o custo da transição entre dois
estados consecutivos.

O algoritmo utiliza uma fila de prioridade,
na qual os nós são ordenados pelo valor
de $g(n)$.

Como os custos das travessias são positivos,
o UCS garante encontrar uma solução de custo
mínimo quando implementado corretamente.

No problema da Ponte e da Tocha, isso significa
encontrar uma sequência de travessias que minimize
o tempo total necessário para transportar todas
as pessoas até a margem final.

\subsection{Busca A*}

A Busca A* é um algoritmo de busca informada que
combina o custo acumulado desde o estado inicial
com uma estimativa do custo restante para alcançar
o estado objetivo.

A seleção dos nós é realizada por meio da função

\[
f(n)=g(n)+h(n),
\]

em que $g(n)$ representa o custo acumulado para
alcançar o nó atual e $h(n)$ representa uma
estimativa do custo necessário para atingir
o estado objetivo.

O algoritmo utiliza uma fila de prioridade
para selecionar o nó com menor valor de $f(n)$.

Dessa forma, o A* considera tanto o custo
do caminho já percorrido quanto uma estimativa
do custo restante, buscando orientar a exploração
para caminhos com menor custo total estimado.

A função heurística $h(n)$ deve ser admissível,
ou seja, não pode superestimar o custo mínimo
real necessário para alcançar o objetivo.

Para solucionar o problema da Ponte e da Tocha, 
propõe-se uma função heurística baseada nos tempos de 
travessia dos indivíduos e na posição atual da tocha. 
Seja $O$ o conjunto de pessoas na margem de origem e $D$ 
o conjunto de pessoas na margem de destino. A função heurística 
$h(n)$ é definida da seguinte forma:

\begin{itemize}
    \item Se a tocha estiver na \textbf{margem de origem}:
    \[
    h(n) = \max_{p \in O} t(p)
    \]
    \item Se a tocha estiver na \textbf{margem de destino} (e $O \neq \emptyset$):
    \[
    h(n) = \max_{p \in O} t(p) + \min_{p \in D} t(p)
    \]
\end{itemize}
Caso o estado seja o objetivo ($O = \emptyset$), $h(n) = 0$.

Essa heurística é \textbf{admissível}. No primeiro caso, 
para que o objetivo seja alcançado, todas as pessoas no 
conjunto $O$ precisam obrigatoriamente atravessar a ponte. 
Como a ponte suporta no máximo duas pessoas e a travessia 
ocorre na velocidade da pessoa mais lenta, o custo mínimo 
para esvaziar a origem será sempre igual ou superior ao 
tempo da pessoa mais lenta que lá está. 

No segundo caso, se a tocha encontra-se no destino e ainda 
existem pessoas na origem, tem-se a garantia de que pelo 
menos uma viagem de retorno precisará ocorrer. O retorno mais 
rápido possível à margem de origem é realizado pela pessoa 
mais rápida presente no destino ($\min_{p \in D} t(p)$). 
Somando-se essa viagem otimista com o limite inferior de 
travessia dos indivíduos restantes na origem ($\max_{p \in O} t(p)$), 
obtém-se um valor que nunca superestima o custo real restante.

Com uma heurística admissível, custos positivos
e tratamento adequado de caminhos alternativos,
o A* garante encontrar uma solução de custo mínimo.

Quando a função heurística assume valor zero
para todos os estados, o A* utiliza apenas o
custo acumulado e passa a apresentar o mesmo
critério de expansão da Busca de Custo Uniforme.

\section{Tarefa 4 -- Comparação dos Métodos de Busca}

Para comparar os quatro métodos implementados (DFS, BFS, UCS e A*),
o programa foi executado sobre o mesmo grafo de estados e as seguintes
métricas foram registradas: o custo da solução encontrada (tempo total
de travessia), o número de nós expandidos durante a busca e o tempo de
processamento até a obtenção da solução. Além disso, todas as buscas
possuem um limite fixo de nós expandidos (no nosso caso, $10\,000$ nós),
após o qual o programa é encerrado. Isso evita buscas excessivamente
longas ou mesmo laços infinitos, o que é especialmente importante no
DFS, que pode mergulhar em ramos muito profundos antes de encontrar
uma solução.

Para o tempo de processamento, cada algoritmo foi executado $1\,000$
vezes e registraram-se a média e o desvio padrão das medições, feitas
com a função \texttt{time.perf\_counter()} da biblioteca padrão do
Python. Os resultados obtidos estão resumidos na tabela a seguir:

\begin{table}[h]
\centering
\begin{tabular}{|l|c|c|c|}
\hline
\textbf{Algoritmo} & \textbf{Custo (min)} & \textbf{Nós expandidos} & \textbf{Tempo médio ($\mu$s)} \\
\hline
DFS & 90 & 13 & $9{,}13 \pm 0{,}85$ \\
BFS & 19 & 26 & $13{,}19 \pm 0{,}59$ \\
UCS & \textbf{17} & 26 & $18{,}53 \pm 2{,}63$ \\
A*  & \textbf{17} & \textbf{17} & $25{,}82 \pm 3{,}61$ \\
\hline
\end{tabular}
\caption{Comparação dos quatro métodos de busca (média de $1\,000$ execuções).}
\end{table}

\subsection{Análise dos Resultados}

\textbf{Custo da solução.} O DFS encontrou a pior solução (90 minutos):
por não considerar nem a profundidade nem o custo das arestas, ele
seguiu o primeiro ramo explorado em profundidade e retornou a primeira
solução encontrada, que continha 10 travessias, muitas delas de alto
custo. O BFS encontrou uma solução de 19 minutos, com o menor número
de transições (5 travessias), mas ainda subótima em tempo. O UCS e o
A* encontraram a solução ótima de 17 minutos.

\textbf{Nós expandidos.} O DFS expandiu apenas 13 nós, mas essa
eficiência é acidental: ele teve a ``sorte'' de tropeçar em uma solução
cedo, sem qualquer garantia sobre sua qualidade. BFS e UCS expandiram
26 nós cada, explorando o espaço de estados de forma sistemática. O A*
expandiu apenas 17 nós, cerca de $35\%$ a menos que o UCS, pois a
heurística admissível orientou a busca em direção ao objetivo, podando
caminhos que, embora parcialmente baratos, não levariam à solução ótima.

\textbf{Tempo de processamento.} Como o grafo do problema possui
apenas 32 estados, o tempo de processamento é da ordem de microssegundos
para todos os algoritmos. Curiosamente, o A* apresentou o maior tempo
médio ($25{,}82\,\mu s$), apesar de expandir menos nós. Isso ocorre
porque, neste problema pequeno, o custo das operações da fila de
prioridade (heap) e do cálculo da função heurística a cada expansão
supera o ganho obtido pela redução do número de nós explorados. De
forma análoga, o DFS, que usa apenas uma pilha simples, foi o mais
rápido ($9{,}13\,\mu s$), ainda que tenha retornado a pior solução.
Em problemas com espaços de estados maiores, espera-se que o número
de nós expandidos seja o fator dominante do tempo total, e a redução
obtida pelo A* (de 26 para 17 nós, cerca de $35\%$ a menos que o UCS)
tende a compensar esse custo adicional por nó.

\subsection{Por que o BFS não garante a solução ótima?}

O BFS garante encontrar a solução com o \textbf{menor número de
transições} (arestas), pois explora os nós em ordem crescente de
profundidade. Entretanto, essa garantia só coincide com a solução de
menor custo quando todas as arestas têm o mesmo custo. No problema da
Ponte e da Tocha, as travessias têm custos diferentes
($t(A)=1$, $t(B)=2$, $t(C)=5$, $t(D)=10$), de modo que um caminho com
menos passos pode ser mais caro que um caminho com mais passos. De
fato, o BFS retornou uma solução de 19 minutos com 5 travessias,
enquanto a solução ótima (17 minutos) também utiliza 5 travessias, mas
com escolhas de pares que minimizam o tempo total -- o BFS, cego aos
custos, não tem como fazer essa distinção: ele retorna o primeiro
estado objetivo que atinge na profundidade mínima, qualquer que seja
seu custo acumulado.

Já o UCS expande sempre o nó de menor custo acumulado $g(n)$. Como
todos os custos das travessias são positivos, quando o estado objetivo
é removido da fila de prioridade pela primeira vez, não existe nenhum
caminho ainda não explorado capaz de alcançá-lo com custo menor -- o
que garante a otimalidade. O A* mantém essa garantia ao expandir o nó
de menor $f(n)=g(n)+h(n)$, desde que a heurística $h(n)$ seja
admissível (nunca superestime o custo real restante), o que é o caso
da heurística utilizada, conforme demonstrado na seção anterior. A
vantagem do A* sobre o UCS é que a heurística direciona a expansão,
reduzindo o número de nós explorados (17 contra 26) sem abrir mão da
solução ótima.

\subsection{Conclusão}

Os resultados confirmam a teoria: DFS e BFS são métodos não informados
que não garantem a solução de menor custo em grafos com arestas de
custos distintos, enquanto UCS e A* garantem a otimalidade. O A*
destaca-se como o melhor método para este problema, pois encontra a
solução ótima (17 minutos) expandindo o menor número de nós, graças
à heurística admissível utilizada. Embora seu tempo médio de
processamento tenha sido ligeiramente superior ao do UCS neste grafo
pequeno -- devido ao custo da fila de prioridade e do cálculo da
heurística --, a economia de nós expandidos tende a se tornar a
vantagem decisiva à medida que o espaço de estados cresce. O limite fixo de nós expandidos
atuou como salvaguarda para impedir buscas excessivamente longas,
especialmente no DFS, que em grafos maiores ou com ciclos poderia
se perder em ramos muito profundos sem encontrar solução.

\end{document}
